\chapter{Analiza porównawcza wyników}

Zestawione wyniki pozwoliły porównać stabilność wyjaśnień generowanych metodami Integrated Gradients oraz Occlusion dla dwóch zbiorów danych: MNIST oraz CIFAR-10. Analiza obejmowała dwa typy perturbacji, czyli szum Gaussa oraz atak FGSM. Oba zakłócenia miały odmienny charakter, dlatego ich wpływ na mapy atrybucji oraz wyniki modelu należało interpretować w różny sposób. Szum Gaussa był perturbacją losową, niezależną od modelu, polegającą na dodaniu przypadkowych zmian do wartości pikseli. FGSM był natomiast perturbacją ukierunkowaną, wyznaczaną na podstawie gradientu funkcji straty względem obrazu wejściowego. Oznaczało to, że atak ten był bezpośrednio związany z mechanizmem decyzyjnym modelu i mógł silniej wpływać zarówno na predykcję, jak i na generowane wyjaśnienia.
Warto podkreślić, że najwyższe wartości parametrów perturbacji, wykorzystywane zarówno w wizualizacjach, jak i w kilku ostatnich punktach analizy ilościowej, pełniły przede wszystkim rolę testu obciążeniowego. Ich celem nie było odtworzenie typowych zakłóceń występujących w danych wejściowych, lecz sprawdzenie, jak model oraz metody wyjaśniania zachowywały się w warunkach silnego zaburzenia obrazu. Z tego względu wyniki dla największych wartości $\sigma$ oraz $\epsilon$ należało interpretować ostrożnie. Były one szczególnie przydatne do oceny granic stabilności metod XAI, ale często dotyczyły już sytuacji, w której zmieniały się nie tylko mapy atrybucji, lecz także pewność modelu lub sama predykcja. Istotne znaczenie miała również różnica pomiędzy analizowanymi zbiorami danych. Obrazy MNIST miały prostą strukturę, zawierały pojedynczą cyfrę na jednorodnym tle, a informacja istotna dla klasyfikacji była skupiona głównie w obrębie kształtu znaku. CIFAR-10 zawierał natomiast obrazy naturalne, kolorowe, z bardziej złożonym tłem i większą liczbą szczegółów wizualnych. W konsekwencji mapy atrybucji dla CIFAR-10 były trudniejsze do interpretacji, ponieważ model mógł wykorzystywać nie tylko fragmenty samego obiektu, ale również elementy jego otoczenia oraz kontekst sceny.

Do oceny stabilności wyjaśnień wykorzystano kilka uzupełniających się miar. Podobieństwo cosinusowe pozwalało ocenić globalną zgodność rozkładu wartości atrybucji pomiędzy mapą dla obrazu oryginalnego i mapą dla obrazu zaburzonego. Ponieważ metryka ta porównywała mapy po spłaszczeniu do postaci wektorów, była mniej wrażliwa na samą skalę wartości. Oznaczało to, że mapa wizualnie osłabiona mogła nadal uzyskać umiarkowane podobieństwo cosinusowe, jeśli zachowywała podobny względny układ wartości. Top-K IoU miało bardziej lokalny i restrykcyjny charakter, ponieważ porównywało zbiory pikseli o najwyższych wartościach atrybucji. Spadek tej metryki oznaczał więc, że po perturbacji zmieniało się położenie regionów uznawanych przez metodę za najważniejsze. Trzecim elementem analizy była metryka Deletion AUC. W przeciwieństwie do podobieństwa cosinusowego i Top-K IoU nie mierzyła ona bezpośrednio podobieństwa dwóch map. Sprawdzała natomiast, jak stopniowe usuwanie pikseli wskazanych przez mapę atrybucji wpływało na pewność modelu względem klasy docelowej. Niższa wartość Deletion AUC mogła oznaczać, że usuwane piksele rzeczywiście były istotne dla predykcji, ponieważ ich zastępowanie wartością bazową szybko obniżało pewność modelu. Interpretacja tej metryki wymagała jednak uwzględnienia średniej pewności modelu oraz liczby zmian predykcji. Jeżeli model już przed rozpoczęciem procedury usuwania miał niską pewność względem klasy docelowej albo zmieniał predykcję, to niska wartość Deletion AUC nie musiała świadczyć wyłącznie o jakości mapy. Mogła częściowo wynikać z tego, że klasyfikacja była już silnie zaburzona przez samą perturbację.

W przypadku zbioru MNIST i szumu Gaussa metoda Integrated Gradients wykazywała stopniowy spadek stabilności wraz ze wzrostem wartości $\sigma$. Podobieństwo cosinusowe zmniejszyło się z $0{,}901$ dla $\sigma = 0{,}010$ do $0{,}598$ dla $\sigma = 0{,}800$, natomiast Top-K IoU spadło z $0{,}764$ do $0{,}400$. Oznaczało to, że ogólny rozkład map atrybucji pozostawał częściowo podobny, ale położenie najbardziej istotnych pikseli zmieniało się coraz wyraźniej. Było to zgodne z charakterem Integrated Gradients, ponieważ metoda ta opierała się na gradientach wyznaczanych wzdłuż ścieżki od obrazu bazowego do obrazu wejściowego. Dodanie szumu wpływało na lokalną strukturę obrazu i  zmieniało wartości gradientów, a w konsekwencji również rozkład atrybucji. Wyniki Deletion AUC dla tego samego eksperymentu pokazały inny aspekt zmian. Dla obrazów niezaburzonych średnia wartość $\overline{AUC}{c}$ wynosiła $0{,}279$, natomiast dla obrazów zaburzonych $\overline{AUC}{p}$ początkowo była wyższa i dla $\sigma = 0{,}010$ wynosiła $0{,}411$. Nie oznaczało to poprawy map po perturbacji. Wyższa wartość Deletion AUC wskazywała raczej, że piksele wskazane przez mapę IG dla obrazu zaburzonego były mniej skuteczne w obniżaniu pewności modelu niż piksele wskazane dla obrazu oryginalnego. Innymi słowy, po niewielkim zaszumieniu mapa mogła nadal wyglądać częściowo podobnie do mapy czystej, ale jej funkcjonalna skuteczność w procedurze deletion była słabsza. Wraz ze wzrostem poziomu szumu $\overline{AUC}_{p}$ zaczęło maleć i dla $\sigma = 0{,}800$ osiągnęło $0{,}104$. Ten końcowy spadek należało jednak interpretować łącznie z pewnością modelu, która dla największego szumu zmniejszyła się do $0{,}703$, oraz z liczbą zmian predykcji równą 24. Przy tak silnym zakłóceniu niska wartość AUC wynikała więc nie tylko z mapy atrybucji, ale także z osłabienia samej klasyfikacji. 

W przypadku Occlusion dla MNIST szum Gaussa miał znacznie słabszy wpływ na metryki podobieństwa. Podobieństwo cosinusowe utrzymywało się na bardzo wysokim poziomie i nawet dla $\sigma = 0{,}800$ wynosiło $0{,}890$, natomiast Top-K IoU spadło do $0{,}624$. Były to wartości wyraźnie wyższe niż dla Integrated Gradients. Wynikało to z regionalnego charakteru Occlusion. Metoda ta oceniała wpływ zasłaniania większych fragmentów obrazu, a nie bezpośrednio lokalne gradienty pojedynczych pikseli. W przypadku prostych obrazów MNIST takie podejście było mniej podatne na drobne, losowe zmiany wartości pikseli, ponieważ najważniejsze regiony nadal pozostawały związane z kształtem cyfry. Deletion AUC również potwierdzało stabilniejszy charakter Occlusion dla szumu Gaussa na MNIST. Dla obrazów czystych $\overline{AUC}{c}$ wynosiło $0{,}184$, a dla obrazów zaburzonych wartości $\overline{AUC}{p}$ przez większość zakresu parametru $\sigma$ pozostawały bardzo blisko tej wartości. Dla $\sigma = 0{,}010$ wynosiły $0{,}185$, a dla $\sigma = 0{,}800$ spadły tylko do $0{,}158$. Oznaczało to, że skuteczność regionów wskazywanych przez Occlusion w procedurze deletion zmieniała się powoli. Było to zgodne z tym, że metoda działała na większych obszarach i nie reagowała tak silnie na pojedyncze zakłócone piksele. Dopiero przy najwyższym poziomie szumu interpretację należało uzupełnić o informację o spadku pewności modelu do $0{,}708$ i 25 zmianach predykcji.

Dla zbioru CIFAR-10 przy szumie Gaussa wyniki Integrated Gradients były częściowo zbliżone do MNIST, jeśli analizowano samo podobieństwo cosinusowe. Spadło ono z $0{,}928$ dla $\sigma = 0{,}010$ do $0{,}529$ dla $\sigma = 0{,}800$. Nie można więc jednoznacznie stwierdzić, że szum Gaussa globalnie wpływał na CIFAR-10 znacznie silniej niż na MNIST. Różnica była jednak wyraźniejsza w metryce Top-K IoU. Dla CIFAR-10 wartość ta spadła z $0{,}586$ do $0{,}114$, podczas gdy dla MNIST przy największym szumie wynosiła $0{,}400$. Oznaczało to, że losowa perturbacja szczególnie silnie zaburzała położenie najbardziej istotnych regionów w obrazach naturalnych, nawet jeśli globalny rozkład atrybucji zachowywał jeszcze umiarkowane podobieństwo. Deletion AUC dla Integrated Gradients na CIFAR-10 pokazało, że dla małych wartości $\sigma$ mapy po perturbacji nadal można było interpretować przy stosunkowo stabilnej klasyfikacji. Średnia wartość $\overline{AUC}{c}$ wynosiła $0{,}187$, a wartości $\overline{AUC}{p}$ dla niewielkiego szumu pozostawały blisko tego poziomu. Dla $\sigma = 0{,}030$ wynosiły nawet $0{,}207$, co oznaczało niewielki wzrost względem obrazów czystych. Podobnie jak w przypadku MNIST nie należało traktować tego jako poprawy wyjaśnień, lecz jako sygnał, że mapa po perturbacji mogła wskazywać piksele mniej skuteczne w szybkim obniżaniu pewności modelu. Przy większych poziomach szumu wartości AUC zaczęły się obniżać, ale równocześnie spadała pewność modelu. Dla $\sigma = 0{,}200$ średnia pewność względem klasy docelowej wynosiła $0{,}590$, a liczba zmian predykcji wzrosła do 17. Dla $\sigma = 0{,}800$ pewność spadła do $0{,}112$, a predykcja zmieniła się w 39 przypadkach. W tych punktach wynik Deletion AUC odzwierciedlał więc nie tylko zmianę jakości map, ale także silne zaburzenie decyzji modelu.

W przypadku Occlusion na CIFAR-10 szum Gaussa również powodował spadek stabilności, ale podobieństwo cosinusowe było wyższe niż dla Integrated Gradients. Dla $\sigma = 0{,}800$ wynosiło $0{,}676$, natomiast Top-K IoU spadło do $0{,}152$. Oznaczało to, że regionalna struktura map Occlusion była bardziej odporna globalnie, ale położenie najważniejszych regionów i tak zmieniało się wyraźnie. W obrazach CIFAR-10 istotne informacje mogły być rozłożone pomiędzy obiektem, tłem i kontekstem sceny, dlatego nawet metoda oparta na zasłanianiu większych obszarów była podatna na przesunięcie lub osłabienie regionów uznawanych za kluczowe. Deletion AUC dla Occlusion na CIFAR-10 wskazywało na inną zależność niż dla MNIST. Dla obrazów czystych $\overline{AUC}_{c}$ wynosiło $0{,}278$, czyli było wyższe niż dla Integrated Gradients. Ponieważ niższa wartość Deletion AUC oznaczała szybszy spadek pewności po usunięciu istotnych pikseli, wynik ten sugerował, że dla czystych obrazów CIFAR-10 mapy Integrated Gradients były bardziej skuteczne w procedurze deletion. Occlusion zachowywała większą stabilność przestrzenną, ale jej szersze, regionalne mapy nie zawsze prowadziły do równie szybkiego spadku pewności modelu. Dla obrazów zaszumionych wartości $\overline{AUC}_{p}$ początkowo pozostawały blisko wartości dla obrazów czystych, a dla $\sigma = 0{,}100$ wzrosły do $0{,}298$. Następnie przy większych poziomach szumu malały, osiągając $0{,}177$ dla $\sigma = 0{,}800$. Tego wzrostu przy średnim poziomie szumu nie należało traktować jako poprawy jakości mapy, lecz jako efekt zmiany sposobu, w jaki mapa Occlusion wskazywała regiony istotne po perturbacji. Przy największych wartościach $\sigma$ interpretacja Deletion AUC była dodatkowo ograniczona przez wyraźny spadek pewności modelu i dużą liczbę zmian predykcji.

Porównanie wyników dla szumu Gaussa pokazało więc, że Occlusion była bardziej stabilna przestrzennie od Integrated Gradients w obu zbiorach danych. Było to szczególnie widoczne w podobieństwie cosinusowym, a w wielu przypadkach także w Top-K IoU. Jednocześnie Deletion AUC pokazało, że większa stabilność przestrzenna nie zawsze oznaczała większą skuteczność wskazywania pikseli krytycznych dla predykcji. Dla MNIST Occlusion uzyskiwała niższe AUC dla obrazów czystych niż IG, natomiast dla CIFAR-10 zależność była odwrotna i korzystniejsze wartości AUC uzyskiwało Integrated Gradients.

W przypadku ataku FGSM zależności były bardziej gwałtowne. Dla MNIST metoda Integrated Gradients traciła stabilność stopniowo wraz ze wzrostem wartości $\epsilon$. Podobieństwo cosinusowe spadło z $0{,}853$ dla $\epsilon = 0{,}010$ do $0{,}541$ dla $\epsilon = 0{,}500$, a Top-K IoU z $0{,}707$ do $0{,}390$. W porównaniu ze szumem Gaussa wartości te były niższe już przy małych poziomach perturbacji, co było zgodne z naturą FGSM. Atak ten nie dodawał losowego szumu, lecz modyfikował obraz w kierunku zwiększającym stratę modelu. Ponieważ Integrated Gradients również opierało się na informacji gradientowej, mapy tej metody były szczególnie podatne na zaburzenia związane z takim typem ataku. Deletion AUC dla Integrated Gradients na MNIST wskazywało na szybki spadek wartości $\overline{AUC}_{p}$ wraz ze wzrostem $\epsilon$. Dla $\epsilon = 0{,}010$ wynosiło ono $0{,}369$, natomiast dla $\epsilon = 0{,}500$ spadło do $0{,}047$. Tego spadku nie należało jednak interpretować jako poprawy jakości map po ataku. Równolegle średnia pewność modelu względem klasy docelowej spadła z $0{,}993$ do $0{,}167$, a liczba zmian predykcji wzrosła z 0 do 83. Oznaczało to, że dla silnego FGSM model często przestawał rozpoznawać klasę docelową jeszcze przed procedurą deletion. Niska wartość AUC była więc związana zarówno ze zmianą mapy atrybucji, jak i z degradacją predykcji.

Dla Occlusion na MNIST metryki podobieństwa przy małych wartościach $\epsilon$ pozostawały bardzo wysokie. Dla $\epsilon = 0{,}010$ podobieństwo cosinusowe wynosiło $0{,}999$, a Top-K IoU $0{,}963$. Przy największej wartości $\epsilon = 0{,}500$ spadły one odpowiednio do $0{,}681$ i $0{,}337$. Oznaczało to, że Occlusion lepiej zachowywała ogólną strukturę map niż Integrated Gradients, ale lokalizacja najważniejszych regionów również ulegała silnej zmianie przy większych poziomach ataku. Było to zrozumiałe, ponieważ Occlusion zależała od odpowiedzi modelu na zasłanianie fragmentów obrazu. Jeżeli atak zmieniał zachowanie modelu, zmieniały się także oceny istotności poszczególnych regionów. Wyniki Deletion AUC dla Occlusion na MNIST uzupełniały tę obserwację. Dla obrazów czystych $\overline{AUC}_{c}$ wynosiło $0{,}184$, czyli było niższe niż w przypadku Integrated Gradients. Wskazywało to, że dla niezaburzonych obrazów MNIST regiony wskazywane przez Occlusion skutecznie obniżały pewność modelu po ich usunięciu. Dla obrazów po ataku $\overline{AUC}_{p}$ spadło z $0{,}179$ dla $\epsilon = 0{,}010$ do $0{,}055$ dla $\epsilon = 0{,}500$. Tego przebiegu nie należało jednak interpretować jako jednoznacznej przewagi nad Integrated Gradients w całym zakresie wartości $\epsilon$. Przy większych wartościach ataku wartości $\overline{AUC}_{p}$ dla Occlusion były nawet nieco wyższe niż dla Integrated Gradients, co wynikało z odmiennego charakteru map generowanych przez obie metody oraz z silnej degradacji predykcji przy większych perturbacjach. W związku z tym niska wartość Deletion AUC przy silnym FGSM nie oznaczała automatycznie lepszej jakości mapy, lecz musiała być interpretowana razem ze spadkiem pewności modelu i rosnącą liczbą zmian predykcji.

Dla CIFAR-10 wpływ FGSM był szczególnie silny. W przypadku Integrated Gradients już dla $\epsilon = 0{,}010$ podobieństwo cosinusowe wynosiło $0{,}762$, a Top-K IoU tylko $0{,}305$. Dla $\epsilon = 0{,}500$ wartości te spadły do $0{,}507$ oraz $0{,}088$. Oznaczało to, że nawet niewielka perturbacja wyraźnie zmieniała położenie najważniejszych obszarów map. Było to zgodne z obserwacjami wizualnymi, w których mapy Integrated Gradients po FGSM dla CIFAR-10 często były silnie osłabione. Umiarkowane wartości podobieństwa cosinusowego nie musiały przeczyć tej obserwacji, ponieważ metryka ta mogła pozostawać względnie wysoka mimo spadku intensywności mapy, jeśli zachowany był częściowy układ wartości. Deletion AUC w tym eksperymencie należało analizować szczególnie ostrożnie. Dla obrazów czystych $\overline{AUC}{c}$ wynosiło około $0{,}185$, natomiast dla obrazów po ataku $\overline{AUC}{p}$ spadało od $0{,}145$ do $0{,}070$. Jednocześnie już dla $\epsilon = 0{,}010$ średnia pewność względem klasy docelowej spadła do $0{,}281$, a liczba zmian predykcji wynosiła 29 na 40 przykładów. Przy $\epsilon = 0{,}075$ zmieniły się wszystkie predykcje. Oznaczało to, że dla CIFAR-10 FGSM bardzo szybko zaburzał nie tylko mapy atrybucji, ale także samą klasyfikację. W takich warunkach niska wartość Deletion AUC mówiła mniej o samej jakości mapy, a więcej o działaniu całego układu: modelu klasyfikacyjnego i metody XAI pod wpływem silnej perturbacji.

Dla Occlusion na CIFAR-10 FGSM również powodował duży spadek stabilności. Podobieństwo cosinusowe zmniejszyło się z $0{,}791$ do $0{,}567$, natomiast Top-K IoU z $0{,}339$ do $0{,}079$. Wartości Top-K IoU były bardzo niskie w całym zakresie parametru $\epsilon$, co oznaczało, że po ataku najważniejsze regiony map Occlusion słabo pokrywały się z regionami wskazywanymi dla obrazów oryginalnych. Chociaż Occlusion nie była metodą gradientową, jej wyniki zależały od reakcji modelu na zasłanianie fragmentów obrazu. Gdy FGSM zmieniał predykcję lub istotnie obniżał pewność względem klasy docelowej, mapa Occlusion również traciła zgodność z mapą dla obrazu czystego. Wyniki Deletion AUC dla Occlusion na CIFAR-10 dodatkowo pokazały niemonotoniczny charakter zachowania modelu przy silnym FGSM. Dla obrazów czystych $\overline{AUC}{c}$ wynosiło $0{,}278$, natomiast dla obrazów po ataku wartości $\overline{AUC}{p}$ były niższe, ale nie malały w sposób monotoniczny. Dla $\epsilon = 0{,}075$ wynosiły $0{,}109$, następnie wzrosły do $0{,}180$ dla $\epsilon = 0{,}150$ i $0{,}190$ dla $\epsilon = 0{,}200$, a później ponownie spadły do $0{,}085$ dla $\epsilon = 0{,}500$. Nie oznaczało to poprawy jakości wyjaśnień w środkowym zakresie wartości $\epsilon$. W tym samym czasie liczba zmian predykcji pozostawała bardzo wysoka, a pewność względem klasy docelowej była niska. 

Porównanie FGSM dla MNIST i CIFAR-10 wskazywało, że w przypadku danych naturalnych wyjaśnienia były znacznie silniej destabilizowane przez perturbację gradientową. Dla MNIST przy najmniejszej wartości $\epsilon$ nie odnotowano zmian predykcji, natomiast dla CIFAR-10 już przy $\epsilon = 0{,}010$ zmieniło się 29 z 40 predykcji. Różnica ta była zgodna z charakterem zbiorów. W MNIST kluczowa informacja była skupiona w prostym kształcie cyfry, natomiast w CIFAR-10 model korzystał z bardziej złożonych cech wizualnych i kontekstu. FGSM, jako atak oparty na gradiencie straty, łatwiej zaburzał takie złożone reprezentacje, co przekładało się również na niższą stabilność map atrybucji.
